\section{Finite-Enclosure Toolkit and Nonzero-Gap Obstructions}
\label{sec:finite-enclosure}

Method~3 has one contract.  Local boundary and radial demands first produce a
point or compact set missed by all six open V roles.  Coverage then forces the
object into the open C triangle.  The case closes either because its least
enclosing equilateral triangle has side at least one or because one forced
point lies beyond the C-triangle exit on the same ray.

\begin{figure}[H]
\centering
\resizebox{.98\linewidth}{!}{\begin{tikzpicture}[
  node distance=6mm and 7mm,
  every node/.style={font=\footnotesize,align=center},
  box/.style={draw=black!55,rounded corners=2pt,fill=CenterBlue!7,
    minimum height=11mm,text width=25mm,inner sep=3pt},
  arr/.style={-{Latex[length=2mm]},line width=.65pt,draw=black!60}
]
\node[box] (class) {normalize\\$N_{\rm gap},N_+,(d,t)$};
\node[box,right=of class] (capacity) {local capacity\\$c_{\max},C_\pm$};
\node[box,right=of capacity] (force) {force a point or\\compact set into $U_C$};
\node[box,right=of force] (terminal) {one named terminal\\$\Lambda\ge1$ or uncovered point};
\node[box,right=of terminal] (dispatch) {close one stable\\routing case ID};
\draw[arr] (class)--(capacity);
\draw[arr] (capacity)--(force);
\draw[arr] (force)--(terminal);
\draw[arr] (terminal)--(dispatch);
\end{tikzpicture}
}
\caption{The public dependency spine for Method~3.  Local capacity and
witness generation are proved once; every combinatorial row invokes one named
terminal.}
\label{fig:finite-enclosure-spine}
\end{figure}

The recurring notation is:
\[
 \Lambda(K),\quad c_{\max}(a,b),\quad C_\pm(a,b),\quad
 D_i=(1-\Gamma_i)V_i,\quad J(p,q),\quad K_{410}.
\]
Here $\Lambda(K)$ is the least enclosing equilateral side, $c_{\max}$ is the
own-ray capacity, $C_\pm$ are the permitted neighboring-ray capacities,
$D_i$ is a type-aware radial witness, $J(p,q)$ is the complementary boundary
gap, and $K_{410}$ is the anisotropic one-gap witness retaining all six
actual radial endpoints.

The chapter is intentionally ordered by logical dependency rather than by
case number: endpoint rules, local capacities, witness generators, universal
terminals, exceptional cases, and finally the case-dispatch table.

\subsection{Endpoint, compactness, and continuity}
\label{subsec:new-endpoint-selector-audit}

The arguments above repeatedly pass between open role traces, closed maximal
reaches, and compact witness sets.  We record the precise endpoint statements
used in those passages.

\begin{lemma}[Open trace endpoint]
\label{lem:new-open-trace-endpoint}
Let $U$ be an open equilateral triangle and $T=\overline U$.  If
$T\cap[V,O]$ is the closed interval from $V$ to a point $P$, then
$P\notin U$.  If $Q$ is a noncentral point of a radial arm and $Q\in U$, then
$U$ has positive-length intersection with that arm.
\end{lemma}

\begin{proof}
The first assertion follows because $P$ is a boundary point of the closed
convex intersection.  For the second, a small Euclidean ball about $Q$ is
contained in $U$; its intersection with the arm is a relative open interval.
\end{proof}

\begin{lemma}[Compact set in an open unit triangle]
\label{lem:new-compact-open-shrink}
If a compact set $K$ is contained in an open unit equilateral triangle, then
\[
 \Lambda(K)<1.
\]
\end{lemma}

\begin{proof}
The distance from $K$ to each of the three side lines is positive.  Move the
three side lines inward by a common smaller distance.  Their intersection is
a closed equilateral triangle of side strictly below one and still contains
$K$.
\end{proof}

\begin{lemma}[Continuity of the enclosure number]
\label{lem:new-lambda-continuity}
For nonempty compact sets $K,L$,
\[
 |\Lambda(K)-\Lambda(L)|
 \le\frac3h d_H(K,L),
\]
where $d_H$ is the Hausdorff distance.
\end{lemma}

\begin{proof}
For every unit normal,
\[
 |h_K(n)-h_L(n)|\le d_H(K,L).
\]
Apply the support formula and take the minimum over orientations in both
directions.
\end{proof}

These lemmas justify all singleton-gap limits and every use of a radial trace
endpoint as a center-forced point.

\subsection{Trace-exact AB envelopes at an actual gap}
\label{subsec:trace-exact-ab-envelopes}

The endpoints of a displayed V-gap are actual maximal reaches, not selected
lower bounds.  If the gap on $e_{i,i+1}$ is
\[
 J_i=X_i([\ell,r]),
 \qquad X_i(t)=V_i+t(V_{i+1}-V_i),
\]
then
\[
 \boxed{\ell=B_i,\qquad r=1-A_{i+1}.}
 \tag{6.24a}
\]
Consequently the two incident open traces are exactly
\[
 U_i\cap e_{i,i+1}=X_i([0,\ell)),
 \qquad
 U_{i+1}\cap e_{i,i+1}=X_i((r,1]),
 \tag{6.24b}
\]
and the whole closed interval $J_i$, including a singleton, is missed by the
six open V roles and therefore lies in $U_C$.

For visualization and local overapproximation it is useful to condition the
ordinary AB-union on this actual endpoint.  If $S$ is a closed unit
equilateral triangle with $V_i\in\interior S$, let $A_i(S),B_i(S)$ be its
actual maximal backward and forward reaches.  Define
\[
 \mathcal E_i^{\rightarrow}(a\mid b)
 =\bigcup\left\{S\cap\mathcal C_i:
 A_i(S)\ge a,\ B_i(S)=b\right\},
 \tag{6.24c}
\]
\[
 \mathcal E_i^{\leftarrow}(a\mid b)
 =\bigcup\left\{S\cap\mathcal C_i:
 A_i(S)=a,\ B_i(S)\ge b\right\},
 \tag{6.24d}
\]
where the union ranges over such closed unit triangles and
$\mathcal C_i$ is the corner cone at $V_i$.  We call these sets the
\emph{trace-exact AB envelopes}.  They are source-conditioned subunions of
the ordinary AB-set, and their relevant closed edge sections are exactly
\[
 \mathcal E_i^{\rightarrow}(a\mid b)\cap e_{i,i+1}
 =X_i([0,b]),
\]
\[
 \mathcal E_{i+1}^{\leftarrow}(1-r\mid b')\cap e_{i,i+1}
 =X_i([r,1]).
 \tag{6.24e}
\]
Thus the actual incident roles are locally contained in
\[
 \mathcal E_i^{\rightarrow}(a_i\mid\ell),
 \qquad
 \mathcal E_{i+1}^{\leftarrow}(1-r\mid b_{i+1}),
 \tag{6.24f}
\]
while the interval between the two exact traces is center-forced.

\begin{proposition}[Trace-exact AB-envelope interface]
\label{prop:readable-trace-exact-ab-interface}
At an actual gap $J_i=X_i([\ell,r])$, the endpoint identities are
$\ell=B_i$ and $r=1-A_{i+1}$.  The source-conditioned envelopes have the
exact edge sections in \emph{(6.24e)}, contain the actual incident roles as
in \emph{(6.24f)}, and satisfy
\[
 \mathcal E_i^{\rightarrow}(a\mid b),\quad
 \mathcal E_i^{\leftarrow}(a\mid b)
 \subseteq \mathcal K_i(a,b).
\]
Consequently the ordinary capacities $c_{\max},C_+,C_-$ remain valid upper
bounds while the full closed interval between the two actual open traces is
center-forced.
\end{proposition}

\begin{proof}
The word ``trace-exact'' is essential.  These sets are not obtained by
intersecting an ordinary AB-union with an affine half-plane through the gap
endpoint: a source triangle whose edge trace stops at that endpoint can cross
such a line elsewhere.  The valid operation restricts the source family to
triangles having the prescribed actual maximal reach and then takes their
union.  The valid operation restricts the source family to triangles having the
prescribed actual maximal reach and then takes their union.  This gives the
edge sections and the inclusions into the ordinary $AB$-sets directly.
The endpoint identities follow from the two actual open incident traces.
The full gap interval is therefore missed by all six V roles and belongs to
$U_C$.


\end{proof}

\input{05a_direct_local_calculus}
\input{05b_neighbor_ray_calculus}
\subsection{High-radial output interfaces}
\label{subsec:high-radial-interfaces}

For a nonsupercritical triangle with incoming reach at least $a$ and radial
reach at least $c$, let
\[
 \mathcal B_c(a)=
 \max\{b:(a,b,c)\text{ is admissible and }a+b\le1\}.
\]
For $1/2<c\le\sqrt3/2$, put
\[
 h(c)=\frac{2c}{1+\sqrt{16c^2-3}},
\]
and for $\sqrt3/2<c<1$ put
\[
 \ell(c)=\frac c2\left(1-\sqrt{4c^2-3}\right),
 \qquad r(c)=c-\ell(c).
\]
Also set
\[
 Q_-(a,c)=\frac{\sqrt{a^2-ac+c^2}}c-a,
\]
\[
 Q_+(a,c)=
 \frac{2(1-a^2)c^2}
 {2ac^2+c+
 \sqrt{(2ac^2+c)^2-4(1-c^2)(1-a^2)c^2}}.
\]

\begin{proposition}[Four direct high-radial outputs]
\label{prop:new-four-direct-outputs}
For $1/2<c\le\sqrt3/2$,
\[
 \mathcal B_c(a)=
 \begin{cases}
  1-a,&0\le a\le1-c,\\
  Q_+(a,c),&1-c<a<h(c),\\
  h(c),&a=h(c),\\
  Q_-(a,c),&h(c)<a<c,\\
  1-a,&c\le a\le1.
 \end{cases}
 \tag{6.66}
\]
For $\sqrt3/2<c<1$,
\[
 \mathcal B_c(a)=
 \begin{cases}
  1-a,&0\le a\le1-c,\\
  Q_+(a,c),&1-c<a\le\ell(c),\\
  \ell(c),&\ell(c)<a<r(c),\\
  Q_-(a,c),&r(c)\le a<c,\\
  1-a,&c\le a\le1.
 \end{cases}
 \tag{6.67}
\]
At $a=\ell(c)$, the selected value is $r(c)$; immediately to the right it is
$\ell(c)$.
\end{proposition}

\begin{proof}
The exact support-cell derivation, including the genuine downward jump at
$a=\ell(c)$, is given in Appendix~\ref{subsec:calc-four-direct-outputs}.
\end{proof}

\begin{lemma}[Quarter radial envelope]
\label{lem:new-quarter-radial-envelope}
If
\[
 0\le h_0\le p,
 \qquad p+h_0<1,
\]
then
\[
 \boxed{c_{\max}(p,h_0)\le1-\frac{h_0}{4}.}
 \tag{6.69}
\]
The inequality is strict when $h_0>0$.
\end{lemma}

\begin{proof}
See Appendix~\ref{subsec:calc-quarter-envelope}, where the $L$ and $T$ cells
are checked separately.
\end{proof}

\begin{lemma}[Vd corner radial margins]
\label{lem:readable-vd-corner-margins}
In the exact Vd1/Vd2 corner form, let $a,b$ be the boundary reaches, let
$t>0$, and put $d=\sqrt{t^2+t+1}$.  The own-radial reach $c_0$ and the far
endpoint $u_+$ on a supported forward adjacent arm satisfy
\[
 c_0=\frac{d-a-tb}{t+1}<1-\min\{a,b\},
 \qquad
 u_+=\frac{d-a-tb-1}{t}<1-b.
\]
The reflected inequalities hold for the backward adjacent arm.
\end{lemma}

\begin{proof}
Since $d<t+1$,
\[
 c_0<1-\frac{a+tb}{t+1}\le1-\min\{a,b\},
 \qquad
 u_+<1-b-\frac at<1-b.
\]
\end{proof}

\subsection{Radial witnesses}

Suppose $T_i=\overline{U_i}$ contains selected boundary anchors at reaches
$a_i,b_i$.  Its exact own-ray capacity is $c_{\max}(a_i,b_i)$.  If the
preceding role has positive support on $r_i$, let its exact neighboring-ray
capacity be $C_+(a_{i-1},b_{i-1})$; define the reflected capacity $C_-$ for
the following role.  Terms not permitted by the actual V type are omitted.
Put
\[
 \Gamma_i=
 \max\{c_{\max}(a_i,b_i),C_+(a_{i-1},b_{i-1}),
 C_-(a_{i+1},b_{i+1})\}
\]
and
\[
 D_i=(1-\Gamma_i)V_i.
 \tag{6.25}
\]

\begin{lemma}[Type-aware radial forcing]
\label{lem:new-type-aware-radial-forcing}
Every $D_i$ is missed by all six open V roles.  Hence, under a hypothetical
cover,
\[
 \boxed{D_i\in U_C.}
\]
\end{lemma}

\begin{proof}
If $D_i\in U_i$, openness increases the own-radial demand beyond
$c_{\max}(a_i,b_i)$.  If an adjacent role contains $D_i$, openness either
creates forbidden positive adjacent support or increases the exact neighboring
capacity beyond $C_+$ or $C_-$.  For the three nonlocal roles, writing
$D_i=dV_i$ gives
\[
 \lVert dV_i-V_{i\pm2}\rVert^2=1+d+d^2>1,
 \qquad
 \lVert dV_i-V_{i+3}\rVert=1+d>1.
\]
A unit equilateral triangle has diameter one.  Thus every V role misses
$D_i$, and coverage forces it into $U_C$.
\end{proof}

\begin{lemma}[Common-pair domination]
\label{lem:new-common-pair-domination}
Let $p,q\ge0$ and $p+q\le1$.  Put
\[
 c_*=c_{\max}(p,q),
 \qquad m=\min\{p,q\}.
\]
Then
\[
 c_*\ge1-m,
 \qquad
 C_+(p,q),C_-(p,q)\le1-m\le c_*.
 \tag{6.26}
\]
Consequently any actual local pair dominating $(p,q)$ has all permitted
own- and neighboring-ray capacities at most $c_*$.
\end{lemma}

\begin{proof}
Assume $q=m\le p$.  The demand $(p,q,1-q)$ has one exact support side of
length $q+\max\{p,1-q\}=1$, so it is admissible; hence $c_*\ge1-q$.
The exact neighboring-ray formula has linear value $1-q$, a plateau no
larger than $1-q$, and a final radical branch no larger than its plateau.
Thus $C_+(p,q)\le1-q$.  Reflection treats the other order and $C_-$.  The
last assertion follows from coordinatewise antitonicity.
\end{proof}

\subsection{Common pairs supplied by actual gaps}

\begin{lemma}[One-gap common boundary pair]
\label{lem:readable-one-gap-common-pair}
Assume there is exactly one V-gap
\[
 J=[B_0,1-A_1]=[\ell,r]\subset e_{0,1}
\]
and every V role on the remaining center-free boundary path is
nonsupercritical.  Put
\[
 p=1-r,\qquad q=\ell.
\]
Then every local boundary pair used in the one-gap argument dominates
$(p,q)$ coordinatewise.
\end{lemma}

\begin{proof}
Choose open-overlap coordinates $x_i$ on the other five edges.  The
nonsupercritical inequalities give
\[
 x_1\le r,
 \qquad x_i\le x_{i-1}\ (2\le i\le5),
 \qquad x_5\ge\ell.
\]
The pair at $T_1$ is $(p,x_1)$; the pair at $T_i$, $2\le i\le5$, is
$(1-x_{i-1},x_i)$; and the pair at $T_0$ is $(1-x_5,q)$.  Each coordinate
is at least the corresponding coordinate of $(p,q)$.
Singleton overlaps follow by the endpoint-continuity lemma.
\end{proof}

\begin{lemma}[Two-gap CE2 common boundary pair]
\label{lem:readable-two-gap-common-pair}
In signed CE2 coordinates, suppose both center intervals contain V-gaps.
Put
\[
 p=W-\alpha,
 \qquad q=R-\delta.
\]
If $T_1,\ldots,T_5$ are nonsupercritical, then every one of these five roles
dominates $(p,q)$ coordinatewise.
\end{lemma}

\begin{proof}
The two incident endpoint roles give $A_1\ge q$ and $B_5\ge p$.
Choose handoffs $x_1,\ldots,x_4$ on the four intervening center-free edges.
Nonsupercriticality gives
\[
 p<x_4<x_3<x_2<x_1<1-q.
\]
Reading the incoming and outgoing pair of each of $T_1,\ldots,T_5$ proves
the coordinatewise domination.
\end{proof}
